\section{Newton selection and moment-shell limits}
\label{sec:newton-selected-kernels}
\index{Bergman kernel!Newton selection}

On a product of projective lines, selecting the monomials whose exponents lie in the dilates of a fixed polytope gives a graded family of section spaces for which the lower bound for the potentials can be proved directly, without a general approximation theorem. The limit of the kernel measures depends on the dimension of the polytope. When this dimension is smaller than that of the manifold, the limit is carried by a family of real tori and differs from the mixed curvature measures of the envelope.

Let $X=(\mathbb P^1)^n$ with $n\geq1$, and let
$L=\mathcal O(k_1,\ldots,k_n)$, where each $k_j$ is a positive integer.
Equip $L$ with its product Fubini--Study metric $h$. On the torus
$T=(\mathbb C^*)^n$, in the standard frame $e$, write
\[
 r_j=\log|z_j|^2,\qquad
 \varphi_0(r)=\sum_jk_j\log(1+\mathrm e^{r_j}),\qquad
 \theta=c_1(L,h)=\ddc\varphi_0.
\]
Put $P=\prod_j[0,k_j]$. Fix a nonempty rational polytope $Q$ contained in the interior of $P$ whose affine span has the form
\[
 \operatorname{aff}Q=b+W,\qquad b\in\mathbb Z^n,
\]
with $W\subseteq\mathbb R^n$ a rational linear subspace. Set $d=\dim Q$ and $\Lambda=W\cap\mathbb Z^n$. Let $\mathrm dq_\Lambda$ be the Lebesgue measure on $b+W$ normalized so that a fundamental domain of the lattice $\Lambda$ has volume one; when $d=0$, it is the unit point mass at $b$. Write $\operatorname{vol}_\Lambda(Q)$ for the volume of $Q$ with respect to $\mathrm dq_\Lambda$. In this section the selected spaces are defined directly by monomials, without a selecting potential: we select exactly the monomials with exponents in $mQ$,
\begin{equation}\label{eq:newton-selected-spaces}
 V_m=\operatorname{span}\{z^\alpha e^m:
                         \alpha\in mQ\cap\mathbb Z^n\}.
\end{equation}
Each of these monomials is a global section of $L^m$, since
$0\leq\alpha_j\leq mk_j$. Since $Q$ is convex, products of selected sections are selected, so these spaces form a graded family.

Choose a smooth real torus-invariant function $w$ on $X$ such that the function
$F=\varphi_0+w$ of $r$ has positive definite Hessian $D^2F(r)$ for every $r\in\mathbb R^n$.
Choose a smooth strictly positive torus-invariant probability measure $\nu$ on $X$. On $T$ we then have
$\mathrm d\nu=a(r)\,\mathrm dr\,\mathrm d\tau$, where $a>0$,
$\int a\,\mathrm dr=1$, and $\mathrm d\tau$ is the Haar probability measure on $(S^1)^n$ in the angular variables.
Give $V_m$ the Hilbert norm $\int_X|s|_{h^m}^2\mathrm e^{-mw}\,\mathrm d\nu$, and for an orthonormal basis put
\[
 \mathcal B_m=\sum_j|s_{m,j}|_{h^m}^2,\qquad
 B_m=\mathrm e^{-mw}\mathcal B_m,\qquad
 U_m=m^{-1}\log\mathcal B_m.
\]
As in \cref{sec:partial-selected-space}, the density $B_m$ includes the continuous weight.

\subsection{The growth constraint and its envelope}

Let $H_Q(r)=\max_{q\in Q}q\cdot r$ be the support function of $Q$, and define
\begin{align*}
 \mathcal C_Q&=\{v\in\PSH(X,\theta):
     v+\varphi_0\leq H_Q+C_v\text{ on }T\text{ for some constant }C_v\},\\
 E_Q&=\left(\sup\{v\in\mathcal C_Q:
                  v\leq w\text{ quasi-everywhere on }X\}\right)^*.
\end{align*}
This growth condition replaces the multiplier-ideal condition of \cref{sec:hermitian-partial-envelopes}; we do not identify $E_Q$ with an ideal envelope.

\begin{lemma}\label{lem:newton-growth-envelope}
The gradient $\nabla F\colon\mathbb R^n\to\operatorname{int}P$ is a smooth
diffeomorphism. With $r(q)=(\nabla F)^{-1}(q)$ and
$F^*(q)=q\cdot r(q)-F(r(q))$, set
\[
 G_Q(r)=\max_{q\in Q}\{q\cdot r-F^*(q)\}.
\]
Then $E_Q=G_Q-\varphi_0$ on $T$, and $E_Q=-\infty$ on $X\setminus T$. Moreover, a global section $s$ of $L^m$ belongs
to $V_m$ exactly when
\begin{equation}\label{eq:newton-entire-growth}
 |s/e^m|^2\leq C_s\mathrm e^{mH_Q(r)}\quad\text{on }T
 \quad\text{for some finite }C_s.
\end{equation}
\end{lemma}
\begin{proof}
In the charts at infinity one sees that every global section of $L^m$ is a polynomial of degree at most $mk_j$ in each variable $z_j$. If $s\in V_m$, the growth bound follows from the triangle inequality, since $|z^\alpha|^2=\mathrm e^{\alpha\cdot r}\leq\mathrm e^{mH_Q(r)}$ for $\alpha\in mQ$. Conversely, averaging the growth bound over the angles gives
$\sum_\alpha|c_\alpha|^2\mathrm e^{\alpha\cdot r}
\leq C_s\mathrm e^{mH_Q(r)}$ for the coefficients $c_\alpha$ of $s$. If some $c_\alpha\neq0$ had $\alpha\notin mQ$, a vector $u$ with $\alpha\cdot u>mH_Q(u)$, which exists by separation, would contradict this bound along the ray $r=tu$, $t\to\infty$. Thus the growth condition controls each monomial separately, even though the monomials of $s$ may cancel pointwise.

For $p\in\operatorname{int}P$, put
$\delta_p=\min_j(p_j,k_j-p_j)>0$. Since $w$ is bounded,
\begin{equation}\label{eq:newton-coercivity}
 F(r)-p\cdot r\geq\delta_p\sum_j|r_j|-\|w\|_\infty.
\end{equation}
By strict convexity, $F-p\cdot r$ therefore has a unique minimizer, at which $\nabla F=p$. Conversely, every gradient of $F$ lies in $P$: compare the supporting hyperplanes of the convex function $F$ with the asymptotic slopes of $\varphi_0$. By the inverse function theorem the image of $\nabla F$ is open, so it is contained in the interior of $P$. This proves that $\nabla F$ is a diffeomorphism onto $\operatorname{int}P$.
On a compact neighborhood of $Q$ in $\operatorname{int}P$, the inverse gradient and the first and second derivatives of $F^*$ are bounded, and the Hessian of $F^*$ is positive definite.

We next show that $G_Q-\varphi_0$ is a competitor. In a product chart in which some of the coordinates are inverted, the local potential of $G_Q-\varphi_0$ is
\[
 \max_{q\in Q}\left\{\sum_j c_j(q)\log|\zeta_j|^2-F^*(q)\right\},
 \qquad c_j(q)=q_j\text{ or }k_j-q_j.
\]
All $c_j$ are bounded below by a fixed positive number. The expression
is locally bounded above, tends uniformly to $-\infty$ at the coordinate
hyperplanes when the other coordinates stay bounded, and is locally
integrable by comparison with one fixed summand. At points where it is finite, a maximizing summand gives the submean inequality for plurisubharmonic functions; at points where it equals $-\infty$, this inequality is automatic. These local expressions glue. Boundedness of $F^*$ on $Q$
gives $G_Q-H_Q=O(1)$, and Fenchel's inequality gives $G_Q\leq F$, that is, $G_Q-\varphi_0\leq w$.
Thus $G_Q-\varphi_0$, extended by $-\infty$ on $X\setminus T$, belongs to $\mathcal C_Q$ and satisfies the obstacle condition.

Now let $v$ be any competitor. The obstacle condition holds everywhere: the exceptional pluripolar set has Lebesgue measure zero, and submean values over shrinking balls carry the inequality with the continuous obstacle to every point. Consider the maximum over the angles,
$M(r)=\max_\tau\{v(\mathrm e^{r/2}\tau)+\varphi_0(r)\}$, where $\mathrm e^{r/2}\tau$ denotes the point with coordinates $\mathrm e^{r_j/2}\tau_j$.
By compactness, $M$ is upper semicontinuous. Rotating a point at which the maximum is attained gives the submean inequality for the invariant function $z\mapsto M(\log|z_1|^2,\ldots,\log|z_n|^2)$, which is therefore plurisubharmonic on $T$. It extends across the affine coordinate hyperplanes by the local upper bound. If $M(r)=-\infty$ for some $r$, successive one-variable maximum principles would force this function to be $-\infty$ on the corresponding polydisk, contradicting $v\not\equiv-\infty$. Thus $M$ is finite. In logarithmic coordinates this invariant function depends only on the real parts, so its positive complex Hessian shows that $M$ is convex.

We have $M\leq F$ and $M\leq H_Q+C_v$. Comparing along rays shows that every supporting slope of $M$ lies in $Q$. At any $r$, choose such a slope
$p$. Since $M\leq F$ gives $M^*\geq F^*$, convex duality yields
$M(r)=p\cdot r-M^*(p)\leq p\cdot r-F^*(p)\leq G_Q(r)$.
Hence $v\leq G_Q-\varphi_0$ on $T$. Submean values over shrinking balls and the uniform decay of $G_Q-\varphi_0$ near $X\setminus T$ show that $v=-\infty$ on $X\setminus T$. Since $G_Q-\varphi_0$, extended by $-\infty$, is already upper semicontinuous, this proves the identity for the envelope after upper regularization.
\end{proof}

\subsection{Norms, lattice growth and the lower bound}

\begin{theorem}\label{thm:newton-potential-lower-bound}
The norms are nondegenerate and satisfy the Bernstein--Markov inequality of \cref{sec:partial-bernstein-markov} for the weight $w$.
Writing $N_m=\dim V_m$, we have
\[
 N_m=m^d\operatorname{vol}_\Lambda(Q)+O(m^{d-1})\quad(d\geq1),
 \qquad N_m=1\quad(d=0).
\]
For all sufficiently large $m$,
\begin{equation}\label{eq:newton-potential-error}
 |U_m-E_Q|\leq\frac{C}{m}
       \left(1+\sum_j|r_j|+\log m\right)\quad\text{on }T.
\end{equation}
Consequently $U_m\to E_Q$ locally uniformly on $T$ and in $L^p$ for every finite $p$ with respect to smooth volume forms on $X$, and the closed positive currents $\theta+\ddc U_m$ converge weakly to $\theta+\ddc E_Q$. Both $U_m$ and $E_Q$ equal $-\infty$ on $X\setminus T$.
\end{theorem}
\begin{proof}
Since the volume form is smooth and positive, the norm is positive on every nonzero
section. Orthogonality of the angular Fourier modes diagonalizes it:
\begin{align*}
 A_m(\alpha)&=\|z^\alpha e^m\|_m^2
   =\int_{\mathbb R^n}\mathrm e^{\alpha\cdot r-mF(r)}a(r)\,\mathrm dr,\\
 \mathcal B_m(r)&=\sum_{\alpha\in mQ\cap\mathbb Z^n}
       \frac{\mathrm e^{\alpha\cdot r-m\varphi_0(r)}}{A_m(\alpha)}.
\end{align*}
For $q=\alpha/m$, the function $F(r)-q\cdot r$ has the unique minimum value
$-F^*(q)$, attained at $r(q)$. Taylor's formula on a ball of radius $m^{-1/2}$ about $r(q)$,
with uniform bounds on a compact neighborhood of $r(Q)$, yields
\begin{equation}\label{eq:newton-laplace-bounds}
 c m^{-n/2}\mathrm e^{mF^*(q)}\leq A_m(\alpha)
                         \leq\mathrm e^{mF^*(q)}.
\end{equation}
The upper bound uses only the minimum and $\int a\,\mathrm dr=1$.

Since $b\in\mathbb Z^n$, we have
$(m\operatorname{aff}Q)\cap\mathbb Z^n=mb+\Lambda$.
Counting fundamental cells of $\Lambda$ in $m(Q-b)$ gives the dimension estimate: the cells meeting the relative boundary of $m(Q-b)$ lie in a boundary layer of fixed width, and there are $O(m^{d-1})$ of them. The same cells give convergence of the lattice Riemann sums on $Q$.
Every $q\in Q$ lies within $O(m^{-1})$ of
$Q_m=Q\cap(b+m^{-1}\Lambda)$: move $q$ by $O(m^{-1})$ towards a fixed relative
interior point so that a ball larger than the covering radius of the lattice fits inside $Q$, and then round to a lattice point inside that ball. This works for all
sufficiently large $m$, and for a singleton it is exact.

Put $G_{Q,m}=\max_{q\in Q_m}(q\cdot r-F^*(q))$. The approximation by the grid
gives $0\leq G_Q-G_{Q,m}\leq C(1+\sum_j|r_j|)/m$.
The formula for $\mathcal B_m$ and \eqref{eq:newton-laplace-bounds} give
\[
 G_{Q,m}-\varphi_0\leq U_m\leq G_{Q,m}-\varphi_0+
     \frac{(n/2)\log m+\log N_m-\log c}{m}.
\]
This proves \eqref{eq:newton-potential-error}. The selected monomials vanish
on every component of $X\setminus T$. The functions $\log|z_j|^2$ have finite
$L^p$ norms for each finite $p$ with respect to smooth volume forms, so \eqref{eq:newton-potential-error} gives the global
convergence, and differentiation in the sense of distributions gives the convergence of the currents.

Fenchel's inequality $\alpha\cdot r-mF(r)\leq mF^*(\alpha/m)$ and \eqref{eq:newton-laplace-bounds} also give
$\sup_X B_m\leq C N_m m^{n/2}$. By the Cauchy--Schwarz inequality in the orthonormal
basis,
$\sup_X|s|_{h^m}\mathrm e^{-mw/2}
\leq C m^{d/2+n/4}\|s\|_m$ for $s\in V_m$.
The finitely many initial degrees can be absorbed in $C$, and a polynomial bound
implies the Bernstein--Markov inequality with any positive exponential loss.
\end{proof}

\Needspace{18\baselineskip}
\subsection{The limit measure and the ambient metric}
\index{Equilibrium measure!moment shell}

For $q\in\operatorname{int}P$ and $\tau\in(S^1)^n$, let $\Psi(q,\tau)=\mathrm e^{r(q)/2}\tau$ be the point of $T$ with coordinates $\mathrm e^{r_j(q)/2}\tau_j$. Thus $\Psi(q,\cdot)$ parametrizes the real torus $\{r=r(q)\}$, and we call the image of $Q\times(S^1)^n$ the \emph{moment shell} of $Q$. Here and below, $\langle S^k\rangle$ denotes the non-pluripolar product of a closed positive $(1,1)$-current $S$; in top degree it is the non-pluripolar Monge--Amp\`ere measure of \cref{thm:arbitrary-extended-ma}.

\begin{theorem}\label{thm:newton-moment-shell-limit}
Against every continuous test function on $X$,
\begin{equation}\label{eq:newton-moment-shell}
 \frac{d!}{m^d}B_m\nu\rightharpoonup
 \mu_Q=\Psi_*(d!\,\mathrm dq_\Lambda\,\mathrm d\tau),
 \qquad \Psi(q,\tau)=\mathrm e^{r(q)/2}\tau.
\end{equation}
The mass of $\mu_Q$ is $d!\operatorname{vol}_\Lambda(Q)>0$. If $d=0$, then $\mu_Q$ is the Haar probability measure on one real $n$-torus. If $d<n$, then $\mu_Q$ is singular with respect to smooth volume forms on $X$, and
\[
 \left\langle(\theta+\ddc E_Q)^{d+1}\right\rangle=0.
\]
If $d=n$, then $\mu_Q=\langle(\theta+\ddc E_Q)^n\rangle$.
\end{theorem}
\begin{proof}
For $\alpha\in mQ\cap\mathbb Z^n$ and $q=\alpha/m$, the normalized monomial $z^\alpha e^m$ gives the probability measure
$\mathrm dP_{m,q}=A_m(\alpha)^{-1}
\mathrm e^{m(q\cdot r-F(r))}a(r)\,\mathrm dr\,\mathrm d\tau$ on $T$.
For every $\eta>0$ there is $c_\eta>0$ such that
$F(r)-q\cdot r+F^*(q)\geq c_\eta$ whenever
$|r-r(q)|\geq\eta$, uniformly for $q\in Q$: on a large compact set this follows from the uniqueness of the minima, and outside it from
\eqref{eq:newton-coercivity}, which is uniform on $Q$. Together with \eqref{eq:newton-laplace-bounds}, this gives
$P_{m,q}\{|r-r(q)|\geq\eta\}\leq C m^{n/2}\mathrm e^{-mc_\eta}$.
Hence, for every continuous function on $X$, its integral against $P_{m,q}$ differs from its average over the torus $\{r=r(q)\}$ by $o(1)$, uniformly in $q\in Q$. The identity
$d!m^{-d}B_m\nu=d!m^{-d}\sum_\alpha P_{m,\alpha/m}$ and the
lattice Riemann sums prove \eqref{eq:newton-moment-shell}. When $d<n$, the measure $\mu_Q$ is carried by the moment shell, which has real dimension $n+d<2n$, so it is singular with respect to smooth volume forms.

For the curvature, note that $\theta+\ddc E_Q=\ddc G_Q$ on $T$. By uniform strong convexity of $F^*$ on $Q$, the slope $q(r)$ maximizing $q\cdot r-F^*(q)$ over $Q$ is unique. Adding the variational inequalities at $r$ and $s$ gives
$(r-s)\cdot(q(r)-q(s))\geq c|q(r)-q(s)|^2$.
Thus $q$ is Lipschitz, and comparison of the maximal values gives
$\nabla G_Q=q$. Hence $G_Q$ is $C^{1,1}$, and its Hessian
has rank at most $d$ almost everywhere, since $q(r)\in b+W$.
In logarithmic coordinates, convex mollifications of $G_Q$ converge in $C^1$,
their Hessians stay bounded and converge in $L^p_{\mathrm{loc}}$ for every finite $p$,
and their minors converge in $L^1_{\mathrm{loc}}$. Induction on the degree in the Bedford--Taylor definition of products of bounded plurisubharmonic functions, together with integration by parts, identifies the powers of $\ddc G_Q$ on $T$
with the measures given by these minors. The non-pluripolar product agrees with them on $T$ and
puts no mass on the pluripolar set $X\setminus T$. The bound on the rank therefore proves the
vanishing assertion.

When $d=n$, a maximizer in the interior of $Q$ satisfies $r=\nabla F^*(q)$, so
$G_Q=F$ on the set where $\nabla F(r)\in\operatorname{int}Q$.
Elsewhere the maximizing slope lies in one of the finitely many facets of $Q$. On the set where the maximizing slope lies in a given facet, its component normal to that facet is constant, so the derivative of this component vanishes at almost every point of the set. Indeed, at a point of differentiability and of density, a nonzero derivative would force a cone of positive density outside the set. Consequently
$\det D^2G_Q=\boldsymbol1_{\{\nabla F(r)\in Q\}}\det D^2F$
almost everywhere. In our normalization,
$(\ddc F)^n=n!\det D^2F(r)\,\mathrm dr\,\mathrm d\tau$.
The change of variables $q=\nabla F(r)$ gives the asserted identity for the top product.
\end{proof}

For any Hermitian form $\omega$ on $X$, closed or not, the map $\Psi$
embeds $\operatorname{relint}Q\times(S^1)^n$ onto a real
$(n+d)$-dimensional submanifold $S_Q=\Psi(\operatorname{relint}Q\times(S^1)^n)$; for $d=0$ this is a single torus.
Let $g_\omega(u,v)=\omega(u,Jv)$ be the associated Riemannian metric, and define the function $\mathcal J_\omega>0$ by
$\Psi^*\mathrm d\operatorname{vol}_{\omega,S_Q}
=\mathcal J_\omega\,\mathrm dq_\Lambda\,\mathrm d\tau$, where $\mathrm d\operatorname{vol}_{\omega,S_Q}$ is the Riemannian volume of $S_Q$ for $g_\omega$.
Since the relative boundary has measure zero,
\[
 \mathrm d\mu_Q=\frac{d!}{\mathcal J_\omega\circ\Psi^{-1}}
                    \,\mathrm d\operatorname{vol}_{\omega,S_Q}.
\]
If $d<n$, the mixed measure
$\langle(\theta+\ddc E_Q)^d\rangle\wedge\omega^{n-d}$, where the zeroth power is interpreted as $1$ when $d=0$, has a locally
bounded density on $T$ and gives $S_Q$ zero mass. Hence it differs from
$\mu_Q$.
The limit $\mu_Q$ does not depend on the invariant probability measure $\nu$, although the norms in each degree do.

\subsection{A rank-one selection with nonclosed sampling}

\begin{example}\label{ex:newton-rational-map-kernels}
On $(\mathbb P^1)^2$ take $L=\mathcal O(4,4)$, $w=0$, and
$Q=\{(2+t,2-t):-1/2\leq t\leq1/2\}$.
Then $\Lambda=\mathbb Z(1,-1)$, $\mathrm dq_\Lambda=\mathrm dt$, and
\[
 V_m=\operatorname{span}\{z_1^{2m+j}z_2^{2m-j}e^m:
                          j\in\mathbb Z,\ |j|\leq m/2\},
 \qquad N_m=2\lfloor m/2\rfloor+1.
\]
The limit $\mu_Q$ has mass one: it is the average, over $-1/2\leq t\leq1/2$, of the Haar probability measures on the real tori
$\{|z_1|^2=(2+t)/(2-t),\ |z_2|^2=(2-t)/(2+t)\}$; indeed, these equations describe the inverse gradient map. The moment shell is a
compact three-dimensional part of the hypersurface $|z_1z_2|=1$.
Writing $c(p)=p\log(p/4)+(4-p)\log((4-p)/4)$, we have
$F^*(2+t,2-t)=2c(2+t)$, which gives the envelope explicitly.

The quotient of two adjacent selected monomials is $z_1/z_2$. In homogeneous coordinates
$z_j=v_j/u_j$, the corresponding map $[v_1u_2:u_1v_2]$ to $\mathbb P^1$ is undefined at
$(0,0)$ and $(\infty,\infty)$. Thus the limit arises from radial
concentration rather than from a globally defined submersion as in \cref{sec:partial-submersion}.

Let $\omega_0$ be the sum of the two Fubini--Study forms, and set
$f=\lambda|z_1|^2/(1+|z_1|^2)$ with a constant $\lambda\ne0$.
Then $\omega=\mathrm e^f\omega_0$ is a torus-invariant Hermitian form,
and $\mathrm d\omega=\mathrm e^f\mathrm df\wedge\omega_0\ne0$
wherever $\mathrm df\ne0$. Sampling with
$\nu=\omega^2/\int_X\omega^2$ gives the same limit of mass one, while the curvature remains $\theta=4\omega_0$.
\end{example}

The assumption $b\in\mathbb Z^n$ excludes a parity obstruction: if $Q=\{b\}$ with $b$ half-integral, then $mQ$ contains no lattice point for odd $m$. Under our assumptions, \cref{thm:newton-potential-lower-bound} shows that $V_m\neq0$ for all sufficiently large $m$. Polytopes touching the boundary of $P$, weights or sampling measures that are not torus-invariant, and general multiplier-ideal families require additional arguments. In particular, the growth condition defining $E_Q$ does not by itself identify the ideal envelope of
\cref{sec:hermitian-partial-envelopes}.
